\documentclass[11pt]{article}
% Shared declarations; root documents always import a concrete semantic module.
\usepackage{business-admin}
\usepackage{amsmath,amssymb}
\usepackage{longtable,tabularx,array,booktabs}
\usepackage{enumitem}
\usepackage{ragged2e}
\newcolumntype{Y}{>{\RaggedRight\arraybackslash}X}
\newcommand{\SuiteMeta}[2]{\def\DocTitle{#1}\def\DocSubtitle{#2}\def\DocVersion{1.0}\def\DocStatus{Educational reference}\def\DocOwner{Small Business Management Suite}\def\DocSummary{Original teaching and practice resource; no real-business approval implied.}}
\newcommand{\SuiteOpen}{\makedoctitle\vspace{-0.7em}\noindent\rule{\linewidth}{0.45pt}\par\medskip}
\newcommand{\SourceNote}[2]{\par\medskip\noindent{\footnotesize\textbf{Reading anchor:} Longenecker, Petty, Palich and Hoy, \emph{Small Business Management: Launching \& Growing Entrepreneurial Ventures}, 20th ed., Ch.~#1, textbook p.~#2. This document is independently authored commentary and exercises, not a reproduction of the chapter.}}
\newcommand{\WritingLines}[1]{\par\smallskip\noindent\rule{\linewidth}{0.35pt}\par\vspace{0.35cm}\noindent\rule{\linewidth}{0.35pt}\par\vspace{0.35cm}\noindent\rule{\linewidth}{0.35pt}\par\medskip}
\newcommand{\FormField}[1]{\par\medskip\noindent\textbf{#1}\par\WritingLines{3}}
\newcommand{\SBMCheckBox}{\(\square\)}
\newcommand{\AssessmentRow}[2]{#1 & #2 & \(\square\) & \(\square\) & \(\square\) & \(\square\) & \(\square\)\\\hline}
\newcommand{\AssessmentHeader}{\noindent\begin{tabular}{|p{0.22\linewidth}|p{0.49\linewidth}|c|c|c|c|c|}\hline\textbf{Criterion} & \textbf{Evidence to inspect} & \textbf{0} & \textbf{1} & \textbf{2} & \textbf{3} & \textbf{4}\\\hline}
\newcommand{\AssessmentEnd}{\end{tabular}\par}
\setlength{\emergencystretch}{3em}

\SuiteMeta{Prerequisite-Based Entrepreneurship Curriculum}{16-module venture decision curriculum}
\begin{document}
\SuiteOpen
\section{Learning contract}
The learner must progress from problem evidence to a feasible business, defend the financial model, demonstrate controllable operations, and explain residual risk. The book order is a useful reading sequence; prerequisites below identify the additional dependencies needed for assessed work.
\section{Entry competencies}
\begin{itemize}[leftmargin=*,itemsep=3pt]
\item Basic spreadsheet arithmetic and percentages.
\item Clear separation of assumptions, primary observations and third-party estimates.
\item Ability to construct a short decision memo and cite evidence.
\item Willingness to seek professional advice for legal, tax, accounting and regulatory decisions.
\end{itemize}

\subsection*{Module 1: chapters 1, 2}
\textbf{Read:} The Entrepreneurial Life; Integrity, Ethics, and Social Entrepreneurship.\par
\textbf{Artifact:} Mission, stakeholder consequences, and customer interviews; Consequence mapping, truthful claims, and escalation owner.\par
\textbf{Gate:} learner must substantiate each claimed outcome and document an unresolved risk.\par
\subsection*{Module 2: chapters 3}
\textbf{Read:} Starting a Small Business.\par
\textbf{Artifact:} Interviews, alternatives, operating capabilities, and stop rules.\par
\textbf{Gate:} learner must substantiate each claimed outcome and document an unresolved risk.\par
\subsection*{Module 3: chapters 4, 5}
\textbf{Read:} Franchises and Independent Contracting in the Gig and Sharing Economies; The Family Business and Buyouts.\par
\textbf{Artifact:} Contract review, actual economics, and classification questions for counsel; Verified books, quality of earnings, and asset/legal due diligence.\par
\textbf{Gate:} learner must substantiate each claimed outcome and document an unresolved risk.\par
\subsection*{Module 4: chapters 6}
\textbf{Read:} Visualizing Your Dream: The Business Model and the Business Plan.\par
\textbf{Artifact:} Evidence for market, operations, team, funding and exit assumptions.\par
\textbf{Gate:} learner must substantiate each claimed outcome and document an unresolved risk.\par
\subsection*{Module 5: chapters 7, 8}
\textbf{Read:} The Marketing Plan; The Organizational Plan: Teams, Legal Structures, Alliances, and Directors.\par
\textbf{Artifact:} Channel performance, acquisition cost, retention and data provenance; Accountable owners, separation of duties, professional review.\par
\textbf{Gate:} learner must substantiate each claimed outcome and document an unresolved risk.\par
\subsection*{Module 6: chapters 9, 10}
\textbf{Read:} The Location Plan; Understanding a Firm’s Financial Statements.\par
\textbf{Artifact:} Whole-life operating cost, constraints and alternatives; Accounting identities, definitions, periods and reconciliation checks.\par
\textbf{Gate:} learner must substantiate each claimed outcome and document an unresolved risk.\par
\subsection*{Module 7: chapters 11}
\textbf{Read:} Forecasting Financial Requirements.\par
\textbf{Artifact:} Base/downside/upside, assumptions, sensitivity, minimum cash.\par
\textbf{Gate:} learner must substantiate each claimed outcome and document an unresolved risk.\par
\subsection*{Module 8: chapters 12, 13}
\textbf{Read:} Financing the Small Business; Planning for the Harvest.\par
\textbf{Artifact:} Debt-service headroom, ownership, collateral, documented assumptions; Transferability, documentation, buyer risk and professional review.\par
\textbf{Gate:} learner must substantiate each claimed outcome and document an unresolved risk.\par
\subsection*{Module 9: chapters 14}
\textbf{Read:} Building Customer Relationships.\par
\textbf{Artifact:} Privacy, response ownership, service recovery and churn evidence.\par
\textbf{Gate:} learner must substantiate each claimed outcome and document an unresolved risk.\par
\subsection*{Module 10: chapters 15, 16}
\textbf{Read:} Product Development and Supply Chain Management; Pricing and Credit Decisions.\par
\textbf{Artifact:} Quality, legal review, inventory, lead time and continuity; Contribution margin, customer impact, collections and defaults.\par
\textbf{Gate:} learner must substantiate each claimed outcome and document an unresolved risk.\par
\subsection*{Module 11: chapters 17}
\textbf{Read:} Promotional Planning.\par
\textbf{Artifact:} Attribution, offer accuracy, acquisition cost and measurement.\par
\textbf{Gate:} learner must substantiate each claimed outcome and document an unresolved risk.\par
\subsection*{Module 12: chapters 18}
\textbf{Read:} Global Opportunities for Small Businesses.\par
\textbf{Artifact:} Compliance review, currency, logistics and contingency.\par
\textbf{Gate:} learner must substantiate each claimed outcome and document an unresolved risk.\par
\subsection*{Module 13: chapters 19, 20}
\textbf{Read:} Professional Management and the Small Business; Managing Human Resources.\par
\textbf{Artifact:} Owner, decision rights, performance review and corrective action; Job relevance, authorization, training evidence and counsel review.\par
\textbf{Gate:} learner must substantiate each claimed outcome and document an unresolved risk.\par
\subsection*{Module 14: chapters 21}
\textbf{Read:} Managing Small Business Operations.\par
\textbf{Artifact:} Process ownership, defects, capacity and supplier SLAs.\par
\textbf{Gate:} learner must substantiate each claimed outcome and document an unresolved risk.\par
\subsection*{Module 15: chapters 22}
\textbf{Read:} Managing the Firm’s Assets.\par
\textbf{Artifact:} Cash visibility, receivables aging, working capital and discount assumptions.\par
\textbf{Gate:} learner must substantiate each claimed outcome and document an unresolved risk.\par
\subsection*{Module 16: chapters 23}
\textbf{Read:} Managing Risk.\par
\textbf{Artifact:} Risk owner, impact, likelihood, triggers and recovery tests.\par
\textbf{Gate:} learner must substantiate each claimed outcome and document an unresolved risk.\par
\section{Integrated assessments}
\begin{enumerate}[leftmargin=*,itemsep=5pt]
\item Gate A / opportunity: validated customer pain and an ethical value proposition (after chapter 5).
\item Gate B / venture design: plan, location/entity assumptions and customer segment (after chapter 9).
\item Gate C / financial readiness: reconciled three statements, cash runway, financing choice and exit hypothesis (after chapter 13).
\item Gate D / go-to-market: proof of repeatable customer value and unit contribution (after chapter 18).
\item Gate E / operating resilience: management cadence, personnel controls, process metrics, cash conversion and risk response (after chapter 23).
\end{enumerate}

\section{Scoring requirements}
Each gate uses the related scorecard: 0=missing, 1=asserted, 2=partially supported, 3=supported, 4=independently corroborated. A gate requires a mean at least 3 and no critical control scored below 2. These are instructional thresholds, not lender or accreditation standards.
\end{document}
